Every day in Oman, patients need medical attention that could be provided remotely if the right tools existed. Children in hospitals need comfort when parents cannot be present. Elderly patients need medication reminders and someone to check on their well-being. Rural communities need access to specialist consultations without long journeys.

The global demand for remote healthcare solutions is growing rapidly, with the telepresence robot market expected to experience significant growth. More importantly, Oman's Vision 2040 calls for healthcare innovation that serves all citizens equally, regardless of location \cite{alhinai2024}. Recent studies have shown that telemedicine adoption in Oman's primary healthcare settings demonstrates both promise and areas for improvement \cite{alsalmani2023}, while infrastructure requirements for effective telemedicine networks continue to be refined \cite{alshishtawy2013}. The aim of the project is to design and implement a smart telepresence robot with autonomous indoor navigation and on-device basic triage dialogue for ward corridors and bedside spaces in hospitals in Oman.